%% 20-related.tex — Related Work

\section{Related Work}
\label{sec:related}

\subsection{EDM Alignment and Ontology Mapping in Cultural Heritage}

Early alignment work on DDB-EDM explored mappings to FaBiO~\cite{peroni2012fabio}:
Tan et al.~\cite{tan2021fabio} showed that the granularity of FaBiO's
subclass hierarchy (over 60 document type classes) cannot be reliably
recovered from DDB metadata without per-provider curation, motivating the
coarser WEMI-level approach adopted in \textsc{mocho}.
Peroni et al.~\cite{peroni2013edm} analysed the structural assumptions
embedded in EDM and identified tensions between EDM's flat model and
work-centric bibliographic models such as FRBR~\cite{ifla1998frbr} and
IFLA LRM~\cite{riva2017iflalrm}.
Patricio et al.~\cite{patricio2025weaving} survey bibliographic ontology
alignment more broadly, documenting the persistent gap between library
standards (RDA, BIBFRAME) and aggregator models (EDM, schema.org).
Our work is complementary: we operationalise alignment as a corpus-driven
pipeline rather than a conceptual mapping exercise.

The DDB's own metadata landscape has been documented by
Schulze~\cite{schulze2020ddb}, who analysed object-describing metadata
conventions across sectors, and by Rühle et al.~\cite{ruehle2014ddbedm},
who describe the DDB-EDM design decisions.
Neither addresses automated alignment to domain ontologies.

\subsection{Automated Ontology Alignment}

OAEI campaigns benchmark automated alignment systems on structured ontology
pairs~\todo{cite OAEI}. Systems such as BERTMap~\cite{he2022bertmap} and
OLaLa~\cite{hertling2023olala} achieve strong results on schema-level
alignment tasks where both source and target are well-specified OWL ontologies.
The EDM-to-domain-ontology problem is structurally different: the source is a
deployed metadata profile whose property semantics are determined by
institutional practice, not by the specification.
Automated tools that match \texttt{dc:type} to a target class on the basis of
lexical similarity will produce a single mapping, missing the sector-specific
divergence that corpus analysis reveals.
This does not make automated tools inapplicable in principle, but it requires
a corpus-first step to partition the alignment problem before any tool is
applied.

Gustavo et al.~\cite{gustavo2018migration} migrate a library catalogue to RDA
Linked Open Data using rule-based mapping, reporting challenges similar to
ours: the source schema (MARC) is well-specified, but the semantics of field
values are institution-specific.
Their work validates the rule-based, domain-knowledge-driven approach we adopt
for DDB-EDM.

%% from mistake in drafting iswc related
%BIBFRAME~\cite{bibframe}, developed by the Library of Congress, provides an RDA-aligned linked data model for bibliographic description.
%The DNB (Deutsche Nationalbibliothek)\todo{cite here!} and OCLC\todo{cite here!} have published linked bibliographic data using BIBFRAME and related vocabularies.
%The EDM-to-RDA/FRBR mapping is not a structural similarity problem: distinguishing FRBR levels requires knowledge of cataloguing conventions; splitting \texttt{edm:Agent} into \texttt{frbr:Person} and \texttt{frbr:CorporateBody} requires external authority file lookups (GND); and the semantics of \texttt{dc:type} are sector-dependent and underdetermined by the EDM specification --- the same property name carries controlled vocabulary in library records, free-text strings in museum records, and structural hierarchy codes in archive records.
%Automated schema-matching tools, as evaluated in OAEI campaigns\todo{cite here!}, presuppose a stable, well-specified source schema and cannot recover implicit, sector-specific semantic conventions absent from the specification itself.
%MOCHO~\cite{tan2026navigating} maps \texttt{edm:ProvidedCHO}\footnote{All CH objects in EDM are instances of this class.} to FRBR~\cite{ifla1998frbr} levels via GND\todo{cite here!} Werk links, aligning DDB-EDM to domain- and modality-specific ontologies (RDA\todo{cite here!}/FRBR, RiC-O\todo{cite here!}, MO\todo{cite here!}, ACO~\cite{ceriani2018audio}, VRA\todo{cite here!}) with DC/DCT properties normalized via RDA's full DCT mapping.
%GeMeA applies MOCHO to the full DDB corpus.

\subsection{Cultural Heritage Knowledge Graphs}

ArCo~\cite{carriero2019arco} is the Italian CH knowledge graph, modelled with
CIDOC-CRM~\cite{doerr2003crm} and covering approximately 820K cultural
objects; it provides a SPARQL endpoint.
The CIDOC-CRM alignment is performed through expert curation rather than
automated corpus analysis, which is feasible at ArCo's scale but not at
DDB scale (65M objects across seven sectors).
Tan et al.~\cite{tan2021ddbkg} report a prior DDB knowledge graph effort
(DDB-KG) using a simpler alignment; the present work advances the alignment
layer to a full three-layer corpus-driven pipeline.

\subsection{Bibliographic Linked Data and Standards}

RDA (Resource Description and Access)~\todo{cite RDA Toolkit} and
BIBFRAME~\todo{cite BIBFRAME} provide linked data models for bibliographic
description.
The DNB has published bibliographic linked data aligned to
RDA~\cite{dnb2022rda}.
\textsc{mocho}~\cite{tan2026mocho} extends these initiatives to heterogeneous
cross-domain EDM metadata, where RDA alone is insufficient because museum
and archive records require VRA~Core and RiC-O respectively.
